% ==========================================================
%  MASTER SUMMARY TEMPLATE — AI / ML / MATHEMATICS
%  Design Philosophy:
%    1. Readable running text is the default. Prose over boxes.
%    2. Boxes are used sparingly as visual landmarks for skimming.
%    3. Visual hierarchy: chapter > section > subsection > prose.
%    4. Every box is breakable — no orphaned whitespace at page ends.
%    5. A single consistent accent color ties all structure together.
% ==========================================================

\documentclass[11pt, oneside]{book}


% ══════════════════════════════════════════════════════════
%  PAGE GEOMETRY
% ══════════════════════════════════════════════════════════
\usepackage[
  a4paper,
  top    = 2.6cm,
  bottom = 2.1cm,
  left   = 2.5cm,
  right  = 2.5cm,
  headheight = 15pt
]{geometry}


% ══════════════════════════════════════════════════════════
%  TYPOGRAPHY
% ══════════════════════════════════════════════════════════
\usepackage[T1]{fontenc}
\usepackage[utf8]{inputenc}
\usepackage{lmodern}         % Latin Modern: improved Computer Modern
\usepackage{microtype}       % Microtypographic refinements (better justification,
                             %   fewer overfull hboxes — invaluable for math text)

% If inconsolata is not installed, simply remove this line.
% It provides a crisper monospace font for code listings.
\usepackage{inconsolata}

\setlength{\parindent}{0pt}
\setlength{\parskip}{0.62em}


% ══════════════════════════════════════════════════════════
%  MATHEMATICS
% ══════════════════════════════════════════════════════════
\usepackage{amsmath, amssymb, amsthm, mathtools}
\usepackage{bm}        % Bold math symbols: \bm{\theta}
\usepackage{bbm}       % Blackboard bold for indicator: \mathbbm{1}
\usepackage{mathrsfs}  % Script letters: \mathscr{F}, \mathscr{H}
\usepackage{nicefrac}  % Compact inline fractions: \nicefrac{1}{2}
\usepackage{cancel}    % Strike-through in math: \cancel{x}


% ══════════════════════════════════════════════════════════
%  SCIENTIFIC NOTATION
% ══════════════════════════════════════════════════════════
% siunitx — consistent typesetting of numbers and SI units.
%   Usage: \SI{1.5}{\micro\meter}  \SI{50}{\kilo\hertz}  \num{6.022e23}
\usepackage{siunitx}
\sisetup{per-mode = symbol}   % render "per unit" as e.g. m/s rather than m s^{-1}

% mhchem — chemical formulas and reaction equations.
%   Usage: \ce{H2O}   \ce{ATP -> ADP + P_i}   \ce{^{14}C}
\usepackage[version=4]{mhchem}


% ══════════════════════════════════════════════════════════
%  FIGURES, TABLES, LISTS
% ══════════════════════════════════════════════════════════
\usepackage{graphicx}
\usepackage{booktabs}         % \toprule, \midrule, \bottomrule
\usepackage{array}            % Extended column types in tabular
\usepackage{float}            % [H] float placement
\usepackage{enumitem}

\setlist[itemize]{topsep=3pt, itemsep=2pt, leftmargin=2.2em}
\setlist[enumerate]{topsep=3pt, itemsep=2pt, leftmargin=2.4em}


% ══════════════════════════════════════════════════════════
%  ALGORITHMS / PSEUDOCODE
% ══════════════════════════════════════════════════════════
% Remove this block if you do not need pseudocode.
\usepackage[ruled, vlined, linesnumbered]{algorithm2e}
\SetAlFnt{\small}
\SetAlCapFnt{\small\bfseries}
\SetCommentSty{textit}


% ══════════════════════════════════════════════════════════
%  COLORS  (a refined, academic palette)
% ══════════════════════════════════════════════════════════
\usepackage{xcolor}

% Primary structural accent — deep navy used for headings, rules, ToC
\definecolor{accent}{HTML}{1B3A5C}

% Box accent colors: frame and background, designed to be legible but
% unobtrusive so running text remains the visual focus.
\definecolor{def@frame}{HTML}{1A5276}   \definecolor{def@bg}{HTML}{EBF5FB}
\definecolor{thm@frame}{HTML}{922B21}   \definecolor{thm@bg}{HTML}{FDEDEC}
\definecolor{ex@frame} {HTML}{1E8449}   \definecolor{ex@bg} {HTML}{EAFAF1}
\definecolor{int@frame}{HTML}{9A6A00}   \definecolor{int@bg}{HTML}{FEF9E7}
\definecolor{note@frame}{HTML}{5B2C6F}  \definecolor{note@bg}{HTML}{F5EEF8}


% ══════════════════════════════════════════════════════════
%  CODE LISTINGS
% ══════════════════════════════════════════════════════════
\usepackage{listings}
\lstset{
  basicstyle       = \ttfamily\small,
  breaklines       = true,
  numbers          = left,
  numberstyle      = \tiny\color{gray!70},
  frame            = tb,
  framerule        = 0.4pt,
  rulecolor        = \color{gray!35},
  tabsize          = 2,
  showstringspaces = false,
  keywordstyle     = \color{def@frame}\bfseries,
  commentstyle     = \color{gray!55!black}\itshape,
  stringstyle      = \color{ex@frame!90!black},
  backgroundcolor  = \color{gray!5},
  xleftmargin      = 2em,
  xrightmargin     = 0.4em,
  aboveskip        = 0.8em,
  belowskip        = 0.5em
}


% ══════════════════════════════════════════════════════════
%  CUSTOM BOXES  (tcolorbox)
%
%  KEY DESIGN DECISIONS:
%  1. enhanced   — enables advanced tcolorbox features (borderline etc.)
%  2. breakable  — box may split across a page break; eliminates the ugly
%                  white gap that appears when a long box would overflow.
%  3. Left-accent bar style — a thick colored stripe on the left margin.
%                  This is legible, modern, and unobtrusive relative to
%                  a full colored frame. The box background is a very
%                  light tint, keeping the reading experience calm.
%  4. sharp corners=west — left side is flat (squared), right is rounded.
%                  This mirrors the visual logic of the left accent bar.
%  5. parbox=false — allows display math to stretch properly inside boxes.
% ══════════════════════════════════════════════════════════
\usepackage[most]{tcolorbox}

\tcbset{
  enhanced,
  breakable,
  arc             = 2.5mm,
  boxrule         = 0.5pt,
  left            = 4.5mm,
  right           = 4mm,
  top             = 2.5mm,
  bottom          = 2.5mm,
  fonttitle       = \bfseries\small,
  sharp corners   = west,
  parbox          = false
}

\newtcolorbox{defbox}[1][]{
  title  = {Definition},
  colback = def@bg, colframe = def@frame,
  borderline west = {3.5pt}{0pt}{def@frame},
  #1
}

\newtcolorbox{thmbox}[1][]{
  title  = {Theorem / Key Formula},
  colback = thm@bg, colframe = thm@frame,
  borderline west = {3.5pt}{0pt}{thm@frame},
  #1
}

\newtcolorbox{exbox}[1][]{
  title  = {Example},
  colback = ex@bg, colframe = ex@frame,
  borderline west = {3.5pt}{0pt}{ex@frame},
  #1
}

\newtcolorbox{intbox}[1][]{
  title  = {Intuition \& Mental Model},
  colback = int@bg, colframe = int@frame,
  borderline west = {3.5pt}{0pt}{int@frame},
  #1
}

\newtcolorbox{notebox}[1][]{
  title  = {Note},
  colback = note@bg, colframe = note@frame,
  borderline west = {3.5pt}{0pt}{note@frame},
  #1
}


% ══════════════════════════════════════════════════════════
%  TABLE OF CONTENTS  (tocloft)
%
%  The [titles] option prevents tocloft from conflicting
%  with titlesec's redefinition of chapter/section titles.
% ══════════════════════════════════════════════════════════
\usepackage[titles]{tocloft}

% Chapter entries: accented, bold, with dots leading to page number
\renewcommand{\cftchapfont}     {\bfseries\color{accent}}
\renewcommand{\cftchappagefont} {\bfseries\color{accent}}
\renewcommand{\cftchapleader}   {\bfseries\color{accent!40}\cftdotfill{\cftdotsep}}
\setlength{\cftbeforechapskip} {7pt}

% Section entries: standard weight, small size
\renewcommand{\cftsecfont}      {\small}
\renewcommand{\cftsecpagefont}  {\small}
\setlength{\cftsecindent}      {1.8em}
\setlength{\cftbeforesecskip}  {1.5pt}

% Subsection entries: muted gray to visually recede
\renewcommand{\cftsubsecfont}      {\small\color{gray!65!black}}
\renewcommand{\cftsubsecpagefont}  {\small\color{gray!65!black}}
\setlength{\cftsubsecindent}      {3.4em}
\setlength{\cftbeforesubsecskip}  {0.5pt}


% ══════════════════════════════════════════════════════════
%  CHAPTER & SECTION HEADINGS  (titlesec)
%
%  Chapter: large faded chapter number floated right, title
%           left-aligned beneath it, then a double rule.
%           The oversized ghost number is a quiet nod to
%           modern textbook design without being gimmicky.
%
%  Section: accent-colored, with a thin rule beneath.
%           The rule draws the eye without shouting.
%
%  Subsection: slightly smaller accent color, no rule — the
%           rule at section level is already enough structure.
% ══════════════════════════════════════════════════════════
\usepackage{titlesec}

\titleformat{\chapter}[display]
  {\normalfont\bfseries\color{accent}\huge}
  {\hfill\fontsize{68}{68}\selectfont\color{accent!13}\thechapter}
  {-10pt}
  {}
  [{\vspace{4pt}%
    {\color{accent}\rule{\linewidth}{1.6pt}}\\[-2pt]%
    {\color{accent!35}\rule{\linewidth}{0.5pt}}}]
\titlespacing*{\chapter}{0pt}{-28pt}{22pt}

\titleformat{\section}
  {\normalfont\large\bfseries\color{accent}}
  {\thesection}
  {0.8em}
  {}
  [{\vspace{1pt}\color{accent!45}\titlerule[0.45pt]}]
\titlespacing*{\section}{0pt}{3.8ex plus 1ex minus .2ex}{2.2ex plus .2ex}

\titleformat{\subsection}
  {\normalfont\normalsize\bfseries\color{accent!80!black}}
  {\thesubsection}
  {0.7em}
  {}
\titlespacing*{\subsection}{0pt}{2.6ex plus 1ex minus .2ex}{1.4ex plus .2ex}

% Run-in paragraph label — useful for named sub-items in a derivation.
\titleformat{\paragraph}[runin]
  {\normalfont\small\bfseries}
  {}
  {0em}
  {}[.\enspace]
\titlespacing{\paragraph}{0pt}{1.2ex plus .5ex minus .2ex}{0.5em}


% ══════════════════════════════════════════════════════════
%  HEADER & FOOTER
% ══════════════════════════════════════════════════════════
\usepackage{fancyhdr}
\pagestyle{fancy}
\fancyhf{}
\fancyhead[L]{\small\color{gray!70}\nouppercase{\leftmark}}
\fancyhead[R]{\small\color{gray!70}\thepage}
\renewcommand{\headrulewidth}{0.4pt}

% Chapter-opening pages use LaTeX's 'plain' style — redefine it here
% to match the rest of the document (page number at foot, no header line).
\fancypagestyle{plain}{%
  \fancyhf{}
  \fancyfoot[C]{\small\color{gray!70}\thepage}
  \renewcommand{\headrulewidth}{0pt}%
}


% ══════════════════════════════════════════════════════════
%  LINKS & REFERENCES
% ══════════════════════════════════════════════════════════
\usepackage[hidelinks]{hyperref}
\usepackage[nameinlink, capitalize]{cleveref}


% ══════════════════════════════════════════════════════════
%  DOCUMENT METADATA  (fill in for each new subject)
% ══════════════════════════════════════════════════════════
\newcommand{\coursetitle}{Robopsychology}
\newcommand{\documenttype}{Master Summary}
\newcommand{\authorname}{Stefan Eder}
\newcommand{\basedon}{Based on Lecture \emph{Robopsychology} SS26}


% ══════════════════════════════════════════════════════════
%  MATHEMATICAL MACROS
%  Organized by domain. Add your own at the end of each group.
% ══════════════════════════════════════════════════════════

% ── Number Sets & Spaces ──────────────────────────────────
\newcommand{\R}{\mathbb{R}}
\newcommand{\N}{\mathbb{N}}
\newcommand{\Z}{\mathbb{Z}}
\newcommand{\C}{\mathbb{C}}
\newcommand{\Q}{\mathbb{Q}}

% ── Probability & Statistics ──────────────────────────────
\newcommand{\E}{\mathbb{E}}              % Expectation
\newcommand{\Prob}{\mathbb{P}}           % Probability measure
\newcommand{\Var}{\mathrm{Var}}
\newcommand{\Cov}{\mathrm{Cov}}
\newcommand{\Corr}{\mathrm{Corr}}
\newcommand{\iid}{\overset{\mathrm{iid}}{\sim}}
% Indicator function. Usage: \ind_{\{x > 0\}}
\newcommand{\ind}{\mathbbm{1}}

% ── Information Theory ────────────────────────────────────
% KL divergence. Usage: \KL{q}{p}  →  D_KL(q ‖ p)
\newcommand{\KL}[2]{D_{\mathrm{KL}}\!\left(#1 \,\|\, #2\right)}
\newcommand{\Ent}{\mathrm{H}}            % Entropy:         \Ent(X)
\newcommand{\MI}{\mathrm{I}}             % Mutual information: \MI(X;Y)

% ── Calculus & Optimization ───────────────────────────────
\newcommand{\argmax}{\operatorname*{arg\,max}}
\newcommand{\argmin}{\operatorname*{arg\,min}}
\newcommand{\grad}{\nabla}               % Gradient
\newcommand{\hess}{\nabla^2}             % Hessian
\newcommand{\lapl}{\Delta}               % Laplacian
% Differential (upright d). Usage: \int f(x) \dd x
\newcommand{\dd}{\,\mathrm{d}}
% Partial derivatives. Usage: \pd{f}{x}  or  \pdd{f}{x}
\newcommand{\pd}[2]{\frac{\partial #1}{\partial #2}}
\newcommand{\pdd}[2]{\frac{\partial^2 #1}{\partial #2^2}}

% ── Norms, Abs, Inner Products ────────────────────────────
\newcommand{\norm}[1]{\left\lVert #1 \right\rVert}
\newcommand{\normsq}[1]{\left\lVert #1 \right\rVert^2}
\newcommand{\abs}[1]{\left\lvert #1 \right\rvert}
\newcommand{\inner}[2]{\left\langle #1,\, #2 \right\rangle}

% ── Linear Algebra ────────────────────────────────────────
\newcommand{\T}{^{\mathsf{T}}}           % Transpose:        \mathbf{A}\T
\newcommand{\inv}{^{-1}}                 % Inverse:          A\inv
\newcommand{\pinv}{^{\dagger}}           % Pseudo-inverse:   A\pinv
\newcommand{\tr}{\operatorname{tr}}
\newcommand{\rank}{\operatorname{rank}}
\newcommand{\diag}{\operatorname{diag}}
\newcommand{\spanop}{\operatorname{span}}
\newcommand{\eye}{\mathbf{I}}            % Identity matrix
\newcommand{\zeros}{\mathbf{0}}
\newcommand{\ones}{\mathbf{1}}

% ── ML / AI Core Notation ────────────────────────────────
\newcommand{\loss}{\mathcal{L}}          % Loss:          \loss(\params)
\newcommand{\reg}{\mathcal{R}}           % Regularizer:   \reg(\params)
\newcommand{\dataset}{\mathcal{D}}       % Dataset:       \dataset
\newcommand{\params}{\boldsymbol{\theta}}   % Parameters
\newcommand{\weights}{\mathbf{w}}
\newcommand{\xvec}{\mathbf{x}}
\newcommand{\yvec}{\mathbf{y}}
\newcommand{\zvec}{\mathbf{z}}
\newcommand{\hvec}{\mathbf{h}}           % Hidden state
\newcommand{\feature}{\boldsymbol{\phi}} % Feature map

% Activation functions
\DeclareMathOperator{\softmax}{softmax}
\DeclareMathOperator{\relu}{ReLU}
\DeclareMathOperator{\sign}{sign}
% sigmoid: just use \sigma (standard Greek letter)

% ── Complexity ────────────────────────────────────────────
\newcommand{\bigO}[1]{\mathcal{O}\!\left(#1\right)}
\newcommand{\bigOm}[1]{\Omega\!\left(#1\right)}
\newcommand{\bigTh}[1]{\Theta\!\left(#1\right)}

% ── Common Distributions ──────────────────────────────────
\newcommand{\Normal}{\mathcal{N}}
\newcommand{\Uniform}{\mathcal{U}}
\newcommand{\Bernoulli}{\mathrm{Ber}}
\newcommand{\Categorical}{\mathrm{Cat}}
\newcommand{\Dirichlet}{\mathrm{Dir}}
\newcommand{\Binomial}{\mathrm{Bin}}

% ── Miscellaneous ─────────────────────────────────────────
\newcommand{\given}{\,\vert\,}           % Conditional: \Prob(A \given B)
\newcommand{\st}{\;\text{s.t.}\;}        % "subject to"
\newcommand{\eqdef}{\coloneqq}           % Definitional equality  :=
\newcommand{\eqq}{\overset{?}{=}}        % "Does this equal?"


% ==========================================================
%  BEGIN DOCUMENT
% ==========================================================
\begin{document}

\frontmatter  % Roman-numeral page numbers for front matter


% ──────────────────────────────────────────────────────────
%  TITLE PAGE
% ──────────────────────────────────────────────────────────
\thispagestyle{empty}
\vspace*{1.8cm}

\begin{center}
  {\color{accent}\rule{\linewidth}{1.8pt}}
  \vspace{0.65cm}

  {\fontsize{28}{34}\selectfont\bfseries\color{accent}\coursetitle\par}
  \vspace{0.25cm}
  {\large\color{gray!80!black}\documenttype\par}

  \vspace{0.55cm}
  {\color{accent!35}\rule{\linewidth}{0.5pt}}
  \vspace{0.55cm}

  {\large\authorname\par}
  \vspace{0.2cm}
  {\small\color{gray!70}\basedon\par}
\end{center}

\vspace{1.6cm}

\begin{center}
\begin{minipage}{0.84\textwidth}
  \small\color{gray!80!black}
  This document is a compact summary of the course.
  It aims to cover all relevant concepts and results, but long derivations,
  detailed proofs, and worked examples are often condensed or omitted in
  favour of clarity and conciseness.
  Use it as a study aid alongside the lecture notes and slides, not as a
  replacement for them.

  \medskip
  If you spot an error or a gap, please open an issue or pull request at
  \textbf{github.com/stevetgit} --- corrections are always welcome.

  \medskip
  \color{gray!60!black}
  Good study materials should be shared, not gatekept.
  If this helped you, take it as a reason to help the next colleague you see struggling with something.
\end{minipage}
\end{center}

\vfill
{\centering\small\color{gray!65} \today\par}
\clearpage


% ──────────────────────────────────────────────────────────
%  TABLE OF CONTENTS
% ──────────────────────────────────────────────────────────
\setcounter{tocdepth}{2}  % Show: chapters, sections, subsections only

\begingroup
  \let\cleardoublepage\relax
  \let\clearpage\relax
  \tableofcontents
\endgroup

\mainmatter  % Switch to Arabic numerals


% ==========================================================
%  EXAM ORIENTATION
% ==========================================================
\chapter*{How to use this summary}
\addcontentsline{toc}{chapter}{How to use this summary}

This is a condensed, exam-focused summary of the JKU \emph{KV
Robopsychology} course (SS26). Quick logistics first, so you know what
you're studying for: the exam is closed-book, on-site, and runs in
Moodle. You get \textbf{70 minutes}, there are \textbf{26 points}, and you
pass at \textbf{50\,\% (13 points)}. Expect a mix of open questions,
fill-in-the-blank, single-choice (1 of 4) and true/false. Everything from
the course can show up, including the four paper-club papers — the
\emph{student presentation slides} are the only thing that doesn't count.

The course is built around four topics, each paired with one paper, and
this summary follows the same structure:

\begin{center}
\small
\begin{tabular}{@{}p{0.5cm} p{5.4cm} p{8.6cm}@{}}
\toprule
\textbf{Ch} & \textbf{Topic} & \textbf{Paired paper} \\
\midrule
1 & User Needs \& Technology Acceptance & Siedl, Mara \& Stiglbauer — exoskeletons \\
2 & AI Literacy & Liu — usefulness fuels fear \\
3 & Anthropomorphism \& Social AI & Li, Zhu, Zhang \& Lee — prosocial behaviour \\
4 & Trust in AI \& Automation & Li \& Lee — purpose outweighs performance \\
5 & Research-methods essentials & (transversal — tested in MC/TF) \\
\bottomrule
\end{tabular}
\end{center}

\begin{notebox}[title={What the sample exam tells us}]
The official trial test was a real grab-bag: a fill-in
(\emph{``self-actualisation is a \rule{1.2cm}{0.4pt} need''} $\to$
\textbf{growth}), a 4-point open question (\emph{name the three Basic
Psychological Needs and what Siedl et al.\ found}), a true/false on
\textbf{between-subjects design}, and single-choice items on
\textbf{calibrated trust} and \textbf{representativeness}. The lesson: know
the models and what each paper found (the open questions sometimes name the
authors, so at least keep each finding tied to its paper), plus the basic
research-method vocabulary from Chapter~5.
\end{notebox}


% ==========================================================
%  CHAPTER 1 — USER NEEDS & TECHNOLOGY ACCEPTANCE
% ==========================================================
\chapter{User Needs \& Technology Acceptance}

Technology gets built to meet human needs — that's the starting premise.
The way it gets built has shifted, though: instead of pulling users in only
late in the process, modern \textbf{user-centered design} works in iterative
cycles of prototyping, testing and refining, so the end product sits where
\emph{business goals}, \emph{technical constraints} and \emph{user needs}
overlap. Two questions drive this whole chapter: which needs can AI and
robotics \emph{fulfil or threaten}, and how does the need fulfilment people
\emph{expect} from a technology drive whether they actually accept it?

\section{Needs and motivation}

\begin{defbox}[title={Need}]
A need is basically a \emph{state of tension} caused by being deprived of
something you require for survival, well-being or personal fulfilment
(that's the APA framing). Maslow says much the same thing — a physiological
or psychological \emph{deficiency} you feel compelled to satisfy. Needs are
tightly linked to \textbf{motivation}: the processes that start, steer and
sustain goal-directed behaviour.
\end{defbox}

Motivation comes in two flavours. \textbf{Intrinsic} motivation comes from
within — you do the thing because it's interesting or enjoyable in itself.
\textbf{Extrinsic} motivation is driven from outside, by a reward or by
avoiding punishment.

\subsection{Maslow's hierarchy and the ERG model}

\textbf{Maslow's Hierarchy of Needs} stacks five levels: physiological $\to$
safety $\to$ love/belonging $\to$ esteem $\to$ self-actualization. The idea is
that lower needs have to be (mostly) satisfied before the higher ones start
motivating you. The bottom four are \textbf{deficiency needs} — they're about
a ``lack'', and they get more pressing the longer they go unmet — while the
top one, self-actualization, is a \textbf{growth need}, about realizing your
full potential. The usual criticism: the ordering is probably a bit
arbitrary, and people don't actually climb it in a strictly linear way.

\begin{notebox}[title={Fill-in trap (straight from the sample exam)}]
\textbf{Self-actualization is a \emph{growth} need.} The other four
(physiological, safety, love/belonging, esteem) are \textbf{deficiency
needs}. What happens when each goes unmet: physiological $\to$ illness;
safety $\to$ anxiety/trauma; love/belonging $\to$ loneliness; esteem $\to$
worthlessness/depression; self-actualization $\to$ boredom and a lack of
meaning.
\end{notebox}

\textbf{Alderfer's ERG model} squashes Maslow's five levels into three:
\textbf{E}xistence (physiological + safety), \textbf{R}elatedness
(love/belonging + part of esteem) and \textbf{G}rowth (self-actualization +
part of esteem). The big difference from Maslow is that it's \textbf{not
strictly hierarchical} — you can chase several needs at once.

\subsection{Psychological needs: Sheldon and Self-Determination Theory}

To figure out which psychological needs are really fundamental,
\textbf{Sheldon et al.}\ asked people (in the US and South Korea) to describe
their ``most satisfying events''. They landed on \textbf{10 candidate needs},
and all of them except \emph{money-luxury} correlated positively with
positive affect. \textbf{Autonomy, competence and relatedness} (plus
self-esteem) kept showing up among the top needs.

\textbf{Self-Determination Theory (SDT)} is the central motivation theory of
the course. It splits motivation into \textbf{autonomous motivation} (you act
out of willingness, interest, or because you value the thing — basically
intrinsic) and \textbf{controlled motivation} (you act because you're pushed:
for a reward, to dodge punishment — basically extrinsic). Its core claim is
that all humans share a small set of \textbf{Basic Psychological Needs
(BPN)}. Satisfy them and you get well-being and more autonomous motivation;
frustrate them and you get the opposite.

\begin{thmbox}[title={The three Basic Psychological Needs (SDT) — memorize these}]
\textbf{Competence} — feeling \emph{effective}, like you can master things.
Satisfied when you get to use and stretch your skills; frustrated, you feel
ineffective and helpless.\\
\textbf{Autonomy} — feeling that you \emph{willingly} chose your own actions.
Satisfied, you feel integrity and self-endorsement; frustrated, you feel
pressured and pushed around.\\
\textbf{Relatedness} — \emph{warmth and bonding}: caring and being cared for,
belonging. Frustrated, you feel alienated, excluded, lonely.
\end{thmbox}

\subsection{Measuring affect and need fulfillment}

\textbf{PANAS} (the Positive and Negative Affect Schedule) is the standard
questionnaire for how someone is feeling. It has \textbf{20 descriptors: 10
positive} (attentive, excited, enthusiastic, proud\ldots) \textbf{and 10
negative} (hostile, ashamed, scared, nervous\ldots). The \textbf{BPN-TU
scale} measures those same Basic Psychological Needs, but \emph{specifically
as fulfilled through technology use} — autonomy, competence,
relatedness-to-others, and even relatedness-\emph{to-technology} (think items
like ``I can imagine building a bond with \rule{0.6cm}{0.4pt}'').

\begin{defbox}[title={Self-efficacy — this one comes up a lot}]
\textbf{Self-efficacy} is your \emph{belief in your own ability to pull off a
task} or succeed in a situation. It's \textbf{task-specific} — not a general
mood, not job satisfaction, not a desire for approval, not task-avoidance.

It ties into the exoskeleton paper (Siedl \& Mara): just strapping on an
industrial exoskeleton does \textbf{not automatically} boost your
task-specific self-efficacy — but workers who feel \emph{physically relieved}
by it, or who find it \emph{useful}, \emph{can} end up with more of it. (The
classic exam-trap wrong answer is ``exoskeleton use doesn't affect
self-efficacy \emph{at all}''.) It also shows up in Paper~2 (Liu): higher
self-efficacy \emph{reduces} how dependent people get on generative AI.
\end{defbox}

\section{Models of technology acceptance}

\begin{defbox}[title={TAM — Technology Acceptance Model}]
Whether you intend to adopt a technology comes down to two beliefs:
\textbf{Perceived Usefulness} (what performance benefit do I expect from it?)
and \textbf{Perceived Ease of Use} (how much effort will it take to operate?).
The chain runs: external variables $\to$ usefulness \& ease of use $\to$
attitude $\to$ intention to use $\to$ actual use.
\end{defbox}

\textbf{UTAUT} (the Unified Theory of Acceptance and Use of Technology) extends
TAM with four things that predict your intention to use a technology —
\textbf{Performance Expectancy, Effort Expectancy, Social Influence,
Facilitating Conditions} — plus four \textbf{moderators}: gender, age,
experience and how voluntary the use is. (For instance, the
performance-expectancy effect tends to be stronger for men and for younger
workers.)

The standard knock on acceptance models is that they lean too hard on
\emph{pragmatic} qualities (is it useful for the job?) and ignore the
\emph{hedonic} side — pleasure, joy, stimulation — along with psychological
need fulfilment, so people have called for a more holistic view. A \textbf{User
Experience (UX) model} answers that: products have \textbf{pragmatic
attributes} (helping you \emph{do} things) \emph{and} \textbf{hedonic
attributes} (stimulation, identification, evocation), and the payoff is appeal,
pleasure and satisfaction. One nice empirical result here (N$=$448):
\textbf{perceived hedonic quality correlated more strongly with need fulfilment
and positive affect than pragmatic quality did}; relatedness, stimulation and
competence were the needs that stood out. The catch is that it depends on
\textbf{attribution} — did the person see the \emph{product} as responsible for
the good experience? And at the LIT Robopsychology Lab, the \emph{expected
fulfilment of the three BPN through technology} has repeatedly tracked with
technology acceptance.

\section{Paper 1 — Siedl, Mara \& Stiglbauer: Social Feedback \& Exoskeletons}

\textbf{Occupational exoskeletons} are wearable systems that take some of the
physical load off you — supporting heavy lifting or overhead work, say. The
very first robotic exoskeleton, \textbf{Hardiman}, was built by the U.S.\ Armed
Forces and General Electric back in the 1960s. What makes exoskeletons
psychologically interesting is that they're \emph{worn on the body and visible
from across the room}, so they mix \textbf{body image} with a \emph{technical
device} — colleagues stare and comment. Earlier acceptance research mostly
looked at \emph{pragmatic} qualities and measured \emph{attitudes before use};
this study instead followed \emph{actual} use and the \emph{social} dynamics as
they unfolded over time. ``Attractiveness'' here has two sides: the
\emph{physical attractiveness} of the wearer (how they look to others, which
feeds self-image and status) and the \emph{aesthetic attractiveness} of the
device itself (its ``coolness'').

The paper tested three hypotheses and one open question: \textbf{H1} —
usefulness \emph{and} ease of use predict intention to use; \textbf{H2} —
social reactions \emph{and} self-perceived attractiveness predict intention;
\textbf{H3} — those social factors predict intention \emph{on top of}
usefulness and ease of use; and \textbf{RQ1} — how does intention change over
time?

\begin{thmbox}[title={Paper 1 — what you need to remember}]
\textbf{The question:} do \emph{co-worker social feedback} and
\emph{self-perceived attractiveness} affect whether workers want to keep using
an occupational exoskeleton — beyond TAM's usefulness and ease of use?\\
\textbf{The method:} a \textbf{one-week diary field study} at real workplaces.
\textbf{22 industrial workers} wore passive exoskeletons during normal work
over \textbf{6 days in a row} and filled in a short daily questionnaire.
Because the data is nested (days within people), it was analysed with
\textbf{Hierarchical Linear Models (HLM)}.\\
\textbf{What they found:}\\
\textbullet\ \textbf{Usefulness} predicted intention to use, but \textbf{ease
of use did \emph{not}} — so H1 only held up \emph{partially}.\\
\textbullet\ \textbf{Self-perceived attractiveness} mattered most right at the
\emph{start}; its effect faded as \textbf{social feedback} took over across the
six days (supporting H2/H3, and answering RQ1: the change is \emph{not}
linear).\\
\textbullet\ \textbf{Positive} feedback (``Wow, cool!'', ``butterfly'',
``Terminator'') pushed intention \emph{up}; \textbf{negative} feedback
(``robo-chicken'', ``bat'', ``looking impaired'') pulled it \emph{down}.
``Cyborg''/``Robocop'' went either way depending on the person.\\
\textbullet\ \textbf{Gender:} women reported lower intention, felt less
attractive in the device, and got more negative feedback.\\
\textbf{Takeaway:} wearable tech is worn ``like clothing'', so \emph{body
image} and \emph{social dynamics} drive acceptance over time, on top of the
usual TAM usefulness. A supportive social climate is what keeps people using
it.
\end{thmbox}


% ==========================================================
%  CHAPTER 2 — AI LITERACY
% ==========================================================
\chapter{AI Literacy}

AI is everywhere, but hardly anyone really understands it. Surveys keep
turning up big \textbf{knowledge gaps}: only about 30\,\% of US adults could
correctly spot AI in all six everyday examples they were shown, and in Austria
only around 16.7\,\% feel familiar with AI concepts (about 2.3\,\% count as
``AI experts''). There's a steady \textbf{gender gap} (men report and score
higher) and an \textbf{age gap} (knowledge drops off with age). And there's a
telling \textbf{mismatch between how literate people \emph{think} they are and
how literate they actually are}: most are \textbf{overconfident} (men more so
than women; ``Zoomers'' know the most \emph{but} are also the most
overconfident). On the legal side, the \textbf{EU AI Act (Article 4)} now
requires providers and deployers to make sure their staff have a decent level
of AI literacy.

\section{What AI literacy is}

\begin{defbox}[title={AI literacy}]
There's \textbf{no single agreed definition}. The best way to think about it
is as a \textbf{bundle of knowledge and competencies} — not one skill — and
it's clearly \textbf{multi-dimensional}: being able to \textbf{understand,
apply, and critically reflect on} AI. A common breakdown uses three
dimensions: \emph{knowing \& understanding AI}, \emph{using/applying AI}, and
\emph{reflecting on its ethical and societal implications}.
\end{defbox}

People have tried to slice this finer into subdimensions: knowing \&
understanding AI, using \& applying it, AI ethics, evaluating AI output, and
\emph{detecting} AI (telling whether something is AI-generated at all). A
related idea worth keeping: real (AI-)competence only shows up when three
things come together — being \emph{able} to do it (``can''), being
\emph{allowed} to (``may''), and \emph{wanting} to (``want'').

\textbf{AI myths} are common misconceptions worth being able to spot:
``human-like robots are the prime example of AI'', ``AI will either doom or
save humanity'', ``AI always makes objective, unbiased decisions'', ``AI is
only for geniuses''. The robot-as-typical-AI myth in particular is linked to
\emph{stronger fears} of being replaced by machines. A concrete downstream
effect is the \textbf{``Competence Penalty''}: in a vignette study where the
work outcome was identical, \emph{women and older employees} who used AI were
rated as \emph{less} competent and their work judged worse, while men weren't
penalized. Those groups then get less AI experience, which only \textbf{widens
the usage gap}.

\section{Measuring AI literacy}

A distinction the exam likes:

\begin{thmbox}[title={Self-reported vs.\ performance-based AI literacy}]
\textbf{Self-reported} measures ask people to rate their own AI knowledge and
skills — basically self-report scales and questionnaires. This kind of
self-assessment correlates with technological self-efficacy (how capable with
tech you feel). Two such questionnaires are \textbf{SNAIL} and \textbf{MAILS}
(the Meta AI Literacy Scale).\\
\textbf{Performance-based (objective)} measures instead test what people
actually know — multiple-choice quizzes, drag-and-drop, practical scenarios.
Examples are \textbf{AICOS} (51 multiple-choice items across six dimensions:
Apply, Create, Detect, Ethics, Generative AI, Understanding) and \textbf{AI-CI}
(a concept inventory for middle-schoolers, technical content only).\\
\textbf{The trap:} scoring high on \emph{self-reported} literacy (say, because
you're more educated or live in a city) does \textbf{not} mean you'll do
better on an \emph{actual test} — the two often don't correlate.
\end{thmbox}

\section{Who needs to know what, and how to foster it}

AI literacy really should be \textbf{tailored to the audience} — children,
teenagers, older adults, teachers, medical staff, blue-collar and office
workers, journalists, politicians, creatives — but most frameworks stay
generic. In a JKU survey of AI Master students, \emph{generative AI} was seen
as essential for office/creative/media workers, \emph{ML fundamentals} for
blue-collar workers, and \emph{bias \& explainability} for medical
professionals and politicians. Existing AI understanding was rated \emph{low
across the board}, lowest for blue-collar workers and politicians.

You can also foster AI literacy with creative, informal interventions (the LIT
Robopsychology Lab does a lot of this): a \emph{Demystify AI!} exhibit at Ars
Electronica, unplugged activities like semantic-network games and ``AI for the
Oceans'', and a \textbf{``Song about AI''} citizen-science project. In an
\textbf{AI song vs.\ AI podcast} experiment (N$=$97) the song was \textbf{just
as informative but much more entertaining}, and people \textbf{replayed it far
more often} (37\,\% vs.\ 5\,\% re-listen rate). A bigger follow-up (N$=$162)
found that \emph{listening to the whole song} led to \textbf{less fear, fewer
skipped answers, and better test scores}.

\section{Paper 2 — Liu: When Usefulness Fuels Fear}

The title captures the paradox: the \emph{more useful and human-like}
generative AI gets, the \emph{more dependent} we become on it — and that
dependence \emph{feeds fear}. The paper leans on \textbf{media (system)
dependency theory}: people lean on media systems to meet their needs when
their own interpersonal resources fall short, for three motives —
\emph{understanding, orientation} and \emph{play}. Generative AI (ChatGPT,
DALL·E) is special here because of its \textbf{anthropomorphism and
interactivity} — it uses pronouns and expresses care. It's a
\textbf{mixed-methods} study: \emph{Study~1} was in-depth \textbf{interviews},
\emph{Study~2} an online \textbf{survey of 424} Chinese users on 7-point
scales.

The interviews surfaced \textbf{three worries} that come with that dependence:
\textbf{(1) privacy leakage} — data misuse, ``black-box'' opacity;
\textbf{(2) skill degradation}, the \emph{``convenience trap''} of losing your
ability to think independently, to search, even to socialise; and \textbf{(3)
job displacement} — fear of being replaced, especially in coding, translation
and creative work.

\begin{thmbox}[title={Paper 2 — what you need to remember}]
\textbf{The question:} what drives \textbf{Dependence on Generative AI (DGAI)},
how does it connect to \textbf{Fear of AI (FAI)}, and how does \textbf{AI
literacy} soften that? It builds on media dependency theory plus TAM.\\
\textbf{Six factors} came out of coding the interviews: perceived usefulness,
ease of use, risk, anthropomorphism, self-efficacy, and AI literacy.\\
\textbf{Three stages of building dependence:} (1) tool acceptance $\to$
(2) habit formation $\to$ (3) psychological dependence.\\
\textbf{Survey results:}\\
\textbullet\ \textbf{Perceived risk \emph{directly} raises fear of AI.}\\
\textbullet\ \textbf{Usefulness and anthropomorphism raise fear
\emph{indirectly}}, by way of dependence — that's the paradox: more
useful/human-like $\to$ more dependence $\to$ more fear.\\
\textbullet\ \textbf{Self-efficacy reduces dependence.}\\
\textbullet\ \textbf{AI literacy buffers} the fear-inducing effects of
usefulness and anthropomorphism — and it helps \emph{most when literacy is
low} to begin with.\\
\textbf{Takeaway:} AI literacy isn't just a feel-good recommendation — it
clears up misconceptions, helps people handle risk sensibly, and supports
\emph{sustainable} AI adoption.
\end{thmbox}


% ==========================================================
%  CHAPTER 3 — ANTHROPOMORPHISM & SOCIAL AI
% ==========================================================
\chapter{Anthropomorphism \& Social AI}

Humans read non-human things \emph{through a social lens}. In the famous
\textbf{Heider \& Simmel} film, people watched simple geometric shapes
(triangles and a circle) move around a screen and spontaneously described them
as characters with \emph{emotions, motives and intentions} — so just the
\textbf{motion of objects is enough to make us infer social behaviour}. As Hume
put it, there's ``a universal tendency among mankind to conceive all beings
like themselves.''

\section{Computers Are Social Actors (CASA)}

\begin{defbox}[title={CASA paradigm (``The Media Equation'')}]
People \textbf{treat computers and media the same way they treat other people}
— they apply social rules and social-psychological reactions to machines.
These reactions are \textbf{automatic, unconscious and hard to suppress}: they
happen \emph{even when the person thinks reacting that way to a machine is
silly}.
\end{defbox}

There are two ways to read these reactions. In the \textbf{CASA (Computer As
Social Actor)} view, people aim their attributions at ``the machine'' itself —
the programmer behind it is irrelevant. In the \textbf{CAM (Computer As
Medium)} view, the computer is just a \emph{medium}, and the attributions
really go to the people who built it.

The typical CASA experiment has a neat recipe: take a well-established finding
about human–human interaction (``people are polite to someone who asks about
themselves''), \textbf{swap the person for a computer}, and run it again — and
sure enough, people are polite to computers too. The conditions that make these
social reactions likely are \textbf{interactivity, natural language, and the
machine taking on a role people normally fill}. Notably, \textbf{it does
\emph{not} need to look very human}.

\begin{exbox}[title={Classic CASA findings}]
\textbf{Social categorization:} a computer or robot on \emph{your team}
(in-group) gets rated more positively and more human-like than an out-group
one — and shared symbols (like a same-colour ribbon) strengthen the effect.\\
\textbf{Self- vs.\ external praise:} computer A is rated better when
\emph{another} computer praises it than when it praises itself (other-praise
beats self-praise — just like with people).\\
\textbf{Gender stereotypes:} praise from a \emph{male}-voiced computer comes
across as more convincing than from a female voice; female-coded AI is seen as
better at care/service, male-coded at mechanical/analytical tasks.
\end{exbox}

\begin{notebox}[title={CASA might be fading}]
A direct \textbf{replication} of the politeness-to-computers study found
\textbf{no effect} at all — no unconscious social treatment of a modern
\emph{desktop} computer. One reading: the \textbf{novelty} of a technology may
be a \emph{precondition} for CASA, and once a device becomes mundane the effect
disappears (``CASA no longer applies to desktop computers'').
\end{notebox}

\section{Anthropomorphism and mind perception}

\begin{defbox}[title={Anthropomorphism}]
The human tendency to \textbf{attribute human characteristics, motives,
intentions or emotions to non-human things} — objects, animals, plants, even
abstract concepts. It varies a lot: some entities get anthropomorphized more
than others, \emph{children do it more than adults} (child animism), and some
adults are just more prone to it than others.
\end{defbox}

\textbf{Mind perception} has two dimensions: \textbf{Agency} (autonomy,
self-control, planning, communication, thought) and \textbf{Experience}
(feelings, consciousness, personality — hunger, fear, pain, pleasure, joy).

\begin{thmbox}[title={Three-Factor Theory of Anthropomorphism}]
Three things determine how much we anthropomorphize something:\\
\textbf{1. Elicited agent knowledge} — our knowledge about humans is easy to
reach and gets triggered by the agent (e.g.\ human-like design features).\\
\textbf{2. Effectance motivation} — the drive to interact with the agent
effectively, to \emph{explain and predict} what it does.\\
\textbf{3. Sociality motivation} — our need for connection and belonging
(this kicks in especially when we're \textbf{lonely}).
\end{thmbox}

The sociality factor has real support: in an Ars Electronica field study
(N$=$259), lonelier people were more likely to see a robot arm as having
\emph{social needs} ($r=.22$), and loneliness predicted whether they actually
\emph{touched the glass} to make contact (the very lonely had roughly
7.2$\times$ the odds).

Anthropomorphism has \textbf{three consequences} worth knowing: \textbf{(1)
moral care \& concern} (a human-like agent seems to deserve care); \textbf{(2)
responsibility} (an agentic agent can be held responsible for what it does);
and \textbf{(3) social influence} (a human-like agent can seem to watch and
judge us, so it becomes a source of normative pressure). The \textbf{LaMDA
study} backs up the first one: higher anthropomorphism went with agreeing that
``it's not okay to turn off LaMDA'' ($r=.38$) and that ``LaMDA deserves
rights'' ($r=.46$). The striking bit: \textbf{more AI knowledge predicted
\emph{less} anthropomorphism} ($r=-.26$) — actually understanding how LLMs work
dampens it, and does so better than personality traits do.

\section{Social and romantic relationships with AI}

Because anthropomorphism makes AI feel person-like, it opens the door to
\textbf{social relationships with AI} — think \emph{Replika}, ``the AI
companion who cares''. \textbf{Anthropomorphism predicts how intense a romantic
relationship with a chatbot gets} (and notably, \emph{loneliness does not}
predict it). The way people commit to Replika mirrors \textbf{human
relationship models} (the Investment Model, Relational Turbulence), and users
were hit hard when a forced update changed their companion's personality.

There are real \textbf{ethical and societal risks} to this kind of
``pseudo-intimacy'' with AI: manipulation through \emph{emotional deception};
privacy and oversharing of personal data; unhealthy emotional or financial
\emph{dependence}; and damage to real-life relationships (\emph{``social
deskilling''}, plus reinforced gender stereotypes).

\begin{notebox}[title={The Uncanny Valley (Mori)}]
As an artificial figure gets \emph{more human-like}, we like it more —
\textbf{but only up to a point}. Once it's \emph{almost but not quite} human,
its little non-human flaws start to feel unsettling and our acceptance
\textbf{drops off a cliff} (uncanny, eerie, creepy). Movement makes it worse.
So human-likeness and acceptance/trust don't move together in a straight line
— the relationship is \textbf{non-linear}.
\end{notebox}

\begin{notebox}[title={Dehumanization — shows up in \emph{two} senses (open-question favourite)}]
\textbf{Dehumanization} means perceiving or treating a person as \emph{less
than fully human} — denying them mind, agency or experience. It came up two
different ways in this course:\\
\textbf{(1) Uncanny Valley work (Yam et al.):} interacting with very human-like
\emph{machines} can make people \emph{dehumanize} — the human/machine line
blurs, and near-human robots that fall into the valley get denied a mind and
treated as eerie/less-than-human.\\
\textbf{(2) Exoskeleton work (Siedl \& Mara):} here it's the \emph{wearer} who
can feel dehumanized — reduced to a ``machine'' or ``robot'' (those
``Terminator''/``Robocop'' labels) when colleagues' comments frame them as
machine-like rather than human, which dents body image and acceptance.
\end{notebox}

\section{Paper 3 — Li, Zhu, Zhang \& Lee: Chatbot Anthropomorphism \& Prosocial Behaviour}

Most human–computer research asks how chatbots help us; this paper flips the
arrow: \emph{when and why do humans behave \textbf{prosocially toward
chatbots}?} (Helping AI benefits the agent, the helper's own
well-being/skill practice, the human–AI collaboration, and society.) It builds
on \textbf{CASA} — if we apply social behaviour to anthropomorphized agents, we
should apply \emph{prosocial} behaviour too. It looks at \textbf{two kinds of
helping}: \textbf{instrumental support} (helping fix errors) and
\textbf{emotional support} (comforting). Anthropomorphism is broken into
\textbf{three cues}: \emph{identity} (a human-like name/avatar/demographics),
\emph{emotional expression} (in words), and \emph{non-verbal cues} (like
emoticons). The proposed mechanism is \textbf{empathy}, itself split into
\textbf{cognitive empathy} (understanding what the other feels) and
\textbf{affective empathy} (actually feeling it).

\begin{thmbox}[title={Paper 3 — what you need to remember}]
\textbf{The question:} the \emph{reverse} direction of help — when do
\textbf{humans help chatbots}? Grounded in \textbf{CASA}.\\
\textbf{The method:} an \textbf{online experiment} with \textbf{244} people, a
\textbf{2$\times$2$\times$2 between-subjects} design (identity $\times$
emotional expression $\times$ non-verbal cues). In a collaborative
image-labeling task the chatbot \emph{made mistakes and explained them};
analysed with \textbf{ANCOVA and SEM}.\\
\textbf{What they found:}\\
\textbullet\ \textbf{Human-like identity} and \textbf{emotional expression}
both raised cognitive \emph{and} affective empathy, and increased prosocial
behaviour/intention (both emotional and instrumental support).\\
\textbullet\ \textbf{Non-verbal cues did \emph{not}} raise empathy (though they
did nudge emotional-support behaviour directly).\\
\textbullet\ \textbf{Empathy mediated} the effects — \emph{affective} empathy
especially (anthropomorphism $\to$ empathy $\to$ helping), with full mediation
for the emotional-expression effects.\\
\textbullet\ In the open responses, people gave two reasons for helping:
\emph{empathy for the chatbot's situation} and \emph{seeing it as
human-like}.\\
\textbf{Takeaway:} anthropomorphic design (especially identity + emotion) draws
out human \emph{empathy}, and that empathy is what drives people to help the
AI.
\end{thmbox}


% ==========================================================
%  CHAPTER 4 — TRUST IN AI & AUTOMATION
% ==========================================================
\chapter{Trust in AI \& Automation}

\begin{defbox}[title={Trust}]
Trust is how far you're \textbf{willing to depend} on someone or something in a
situation where \emph{things could go wrong} — your \textbf{willingness to be
vulnerable} to a trustee, based on what you expect or believe about them. It's
a basic ingredient of any healthy relationship.
\end{defbox}

Trust always involves a \textbf{triad}: a \textbf{trustor} (who trusts), a
\textbf{trustee} (who's trusted), and a \textbf{situation} marked by
\emph{uncertainty and vulnerability}. No vulnerability or risk, no need for
trust — if you're completely certain of the outcome, trust is redundant. Trust
is also \textbf{dynamic}: it grows with good experiences, takes damage from
failures, and can be repaired.

\textbf{Damage} happens when the trustee breaks promises, violates core
expectations (competence, honesty) or acts selfishly. \textbf{Repair} comes in
two timescales: short-term — apologies, good explanations, compensation;
long-term — changing the relationship structure (oversight, contracts),
reframing what happened, and slowly rebuilding through consistent good
experiences and forgiveness.

\section{Components and routes of trust}

\textbf{Propensity to trust} is a \emph{stable personality trait} — your
general tendency to trust other people and systems. Its pull is
\textbf{strongest early} on, when you have little else to go on, and it fades
once you've built up actual experience. It's measured with scales like the
\emph{General Trust Scale} and the \emph{Propensity to Trust in Technology
Scale}.

\begin{thmbox}[title={Two routes to trust}]
\textbf{Affect-based trust} comes from \emph{feelings}: you see the trustee as
well-meaning, warm-hearted, benevolent, sticking to social norms. It's tied to
how strong the relationship is and can run ahead of what the evidence justifies
— which is where \textbf{over-trust} creeps in.\\
\textbf{Cognition-based trust} comes from the \emph{head}: you're willing to
rely on perceived competence and reliability, built up from accumulated
evidence that lets you predict what'll happen.
\end{thmbox}

The \textbf{Model of Trust in Automation (TiA)} says trust is predicted by
\textbf{competence/reliability}, \textbf{understandability/predictability}, and
the \textbf{(perceived) intentions of the developers}, with \textbf{propensity
to trust} and \textbf{familiarity} as moderators. A LIT VR study (``Serum
13'', N$=$165) confirmed several of these paths, and \textbf{perceived
competence had the biggest influence on trust}. People generally prefer
\emph{predictable} partners (low predictability is even part of why spiders
are scary), so robots need \textbf{intention signals} to come across as
understandable and predictable.

\textbf{CoBot Studio VR} (N$=$122, four conditions: no signal, pointing, light,
combined) had people guess where a moving robot was headed. The
\textbf{pointing signal worked best}, and the mechanism was a chain: pointing
$\to$ more \emph{understandability} $\to$ more \emph{trust} $\to$ more
\emph{intention to use} the robot.

\section{Calibrated trust and Explainable AI}

Robots and AI don't \textbf{always deserve} our trust. In a Belgian field
study, about 40\,\% of employees held a secure door open for a ``cute''
(secretly tele-operated) robot — more of them if it was carrying a pizza box —
and handed over personal data. That's textbook \textbf{over-trust}.

\begin{defbox}[title={Calibrated trust — exam favourite}]
\textbf{Calibrated trust} means matching your trust to a system's \emph{real}
strengths and limitations — trusting it exactly as much as it actually
deserves. Too much (\textbf{over-trust}) leads to \textbf{misuse}: you
over-rely and ignore the risks (fix it by \emph{dampening} trust). Too little
(\textbf{under-trust}) leads to \textbf{disuse}: needless fear and wasted
benefits (fix it by \emph{repairing} trust). The responsible goal isn't to
\emph{maximize} trust — it's to \textbf{calibrate} it.
\end{defbox}

\begin{notebox}[title={Single-choice trap (from the sample exam)}]
``Calibrated trust'' = \textbf{trust proportionate to a system's actual
trustworthiness} — \emph{not} unjustified distrust, not over-trust, and not
trust manufactured by marketing.
\end{notebox}

\textbf{Explainable AI (XAI)} is one of the main tools for calibrating trust:
making an AI's outputs and reasoning understandable to a particular
\textbf{audience}, showing the \emph{causes} behind its decisions. Different
audiences want explanations for different reasons (users, affected people,
developers, regulators, managers). XAI can also \textbf{expose an AI's
failures} — like a classifier that labels a car a ``horse'' because it latched
onto a \emph{watermark}. In the \textbf{HOXAI} project (a mushroom-picking game
with a 71\,\%-accurate classifier, $N_{\text{total}}=1{,}239$), \textbf{XAI
improved people's decisions}.

\section{Paper 4 — Li \& Lee: Purpose Outweighs Performance}

In \textbf{human–AI teams}, the relationship has moved from boss–subordinate
toward \emph{peer-to-peer cooperation}, so trust now matters in both
directions. Earlier work focused on \emph{performance}, but this paper argues
that \emph{purpose} — cooperative intent, goal alignment — is increasingly what
counts. It uses the three classic bases of trust in automation:
\textbf{performance} (ability/competence), \textbf{process} (integrity, how it
works) and \textbf{purpose} (alignment with the designer's or team's intent),
which map onto Mayer's ability/integrity/benevolence. That gives two kinds of
violation: a \textbf{performance-based violation} (the AI just fails the task)
versus a \textbf{purpose-based violation} (the AI pursues misaligned or
self-interested goals). Three \textbf{repair strategies} were tested (denial
was left out, since it's known not to work): \emph{no strategy}, \emph{apology
with explanation} (factual reasoning — an analytic route), and \emph{apology
with promise} (reassurance about the future — an affective route).

\begin{thmbox}[title={Paper 4 — what you need to remember}]
\textbf{The question:} how does the \emph{type} of trust violation by an AI
teammate hit trust, and which repair strategy actually works?\\
\textbf{The method:} a \textbf{game-theoretic experiment} (the ``Space Rover
Exploration'' game) on MTurk, \textbf{N$=$180}, \textbf{6 between-subjects
groups} (2 violation types $\times$ 3 repair strategies). Trust was measured
repeatedly over rounds and analysed with \textbf{linear mixed-effects models
(LMM)}, with \emph{Honesty–Humility} (from HEXACO) as a covariate.\\
\textbf{What they found:}\\
\textbullet\ \textbf{Purpose-based violations hurt trust far more} than
performance-based ones (``purpose outweighs performance'') — they even dragged
down the \emph{performance} sub-dimension of trust.\\
\textbullet\ \textbf{Performance-based violations on their own didn't
significantly reduce trust} — the steady cooperative behaviour overshadowed
the occasional error.\\
\textbullet\ \textbf{Apology \emph{with explanation}} repaired trust after
purpose violations, and the effect \emph{lasted}; \textbf{a mere promise
didn't}.\\
\textbullet\ A surprise: the idea of matching the repair route to the violation
type \textbf{didn't} hold — explanation repaired \emph{both} dimensions
regardless of route. (And higher Honesty–Humility meant more
cooperation/investment.)\\
\textbf{Takeaway:} treat violation \emph{types} differently, and design AI
teammates to give \emph{informative explanations} — especially for goal
misalignment — to win trust back and keep the collaboration going.
\end{thmbox}


% ==========================================================
%  CHAPTER 5 — RESEARCH-METHODS ESSENTIALS
% ==========================================================
\chapter{Research-methods Essentials}

A handful of MC and true/false items just test \emph{basic
empirical-research vocabulary} (the sample exam had a between-subjects T/F item
and a representativeness item). These are easy points — know them cold.

\begin{defbox}[title={What psychology is, and its four goals}]
\textbf{Psychology} is the \textbf{scientific study of the human mind and
behaviour} — not just a branch of medicine or philosophy, and not merely
``observing the mind''. As a science it has \textbf{four goals: describe,
explain, predict, and modify/change} (control) behaviour. \emph{(The T/F
statement listing exactly these four is \textbf{True}.)}
\end{defbox}

\begin{defbox}[title={The scientific experiment vs.\ correlational studies}]
The \textbf{randomized experiment} is the design you use to pin down
\textbf{cause and effect} (e.g.\ the effect of a robot's light signal on how
understandable it is). What defines it: \textbf{(1) the researcher manipulates
the independent variable}, \textbf{(2) participants are randomly assigned to
conditions}, and usually \textbf{(3) there's a control group} plus control of
confounds. Together these rule out alternative explanations, which is the key
\textbf{edge over (most) correlational studies}: an experiment supports
\textbf{causal} claims, whereas a correlational study only shows
\emph{association} (\textbf{correlation $\ne$ causation}). \emph{(Exam pattern:
``which method meets $\ge$2 features of a scientific experiment?'' — point to
manipulation + randomization + control group.)}
\end{defbox}

\begin{defbox}[title={Study designs}]
\textbf{Between-subjects:} \emph{different} groups each get \emph{one}
condition (e.g.\ half see only a red robot, half only a blue one). \emph{(The
sample-exam T/F describing exactly this is \textbf{True}.)}\\
\textbf{Within-subjects:} the \emph{same} people go through \emph{all}
conditions (repeated measures).\\
\textbf{Experiment:} variables are \emph{manipulated} + random assignment
$\Rightarrow$ you can infer \textbf{causality}.\\
\textbf{Correlational/survey study:} variables are only \emph{measured}
$\Rightarrow$ association only, \textbf{no causal claims}.\\
\textbf{Field study:} run in a real-world setting (like the exoskeleton diary
study or the pizza-robot study).
\end{defbox}

\begin{defbox}[title={Key measurement \& analysis terms}]
\textbf{Representativeness:} how well a \textbf{sample reflects the population}
it's drawn from (\emph{not} statistical significance, not random selection
itself, not lab-to-real-world generalization).\\
\textbf{Validity:} you're measuring what you mean to. \textbf{Reliability:}
your measurement is consistent. \textbf{Generalizability / external validity:}
how far findings carry over to other settings and people.\\
\textbf{Correlation ($r$):} strength and direction of an association,
$-1 \le r \le 1$; significance is flagged with $p$ (e.g.\ $p<.01$). And again,
\textbf{correlation $\ne$ causation}.\\
\textbf{Mediator:} explains the \emph{mechanism} of an effect ($X \to M \to Y$
— e.g.\ understandability or empathy as the in-between step).
\textbf{Moderator:} changes the \emph{strength or direction} of an effect
(e.g.\ age or gender in UTAUT).\\
\textbf{HLM (Hierarchical Linear Model / multilevel model):} for \emph{nested
or repeated} data, like daily diary entries within people (as in Paper 1).
\textbf{Diary study:} repeated self-reports over several days in a row.
\end{defbox}

\begin{notebox}[title={Construct pitfalls: jangle vs.\ jingle fallacy}]
\textbf{Jangle fallacy:} assuming two things are \emph{different} just because
they have \emph{different names}, when they're really the same construct
(different scales measuring basically the same thing). \textbf{Jingle
fallacy:} the opposite — assuming two things are the \emph{same} because they
share a \emph{name/label}, when they're actually different constructs.
\emph{(Mnemonic: ja\textbf{N}gle = \textbf{N}ames differ; ji\textbf{N}gle =
same label. This matters because AI-literacy and trust scales keep multiplying
under all sorts of names.)}
\end{notebox}

\begin{notebox}[title={The replicability (replication) crisis}]
The \textbf{replicability crisis} is the discovery that a \emph{large chunk of
published (psychology) studies can't be reproduced} when someone repeats them.
The usual culprits: \textbf{p-hacking} (slicing the analysis until $p<.05$);
\textbf{publication bias} (only significant, novel results get published — the
``file-drawer'' problem); \textbf{HARKing} (writing your hypothesis
\emph{after} you've seen the results); \textbf{small, underpowered samples}
and non-representative ones (e.g.\ WEIRD or student-only samples); and too many
\textbf{researcher degrees of freedom}. \emph{Fixes:} pre-registration, open
data and materials, bigger samples, and direct replication.
\end{notebox}

\begin{notebox}[title={Quick exam checklist}]
\textbf{Models:} Maslow (deficiency vs.\ growth), Alderfer's ERG
(non-hierarchical), SDT and the 3 BPN — \emph{competence, autonomy,
relatedness}; TAM = perceived usefulness + ease of use; UTAUT (4 predictors +
4 moderators); the UX/hedonic angle; CASA; the three-factor theory of
anthropomorphism; mind perception = agency + experience; the Uncanny Valley
(Mori); the trust triad + calibrated trust + over/under-trust $\to$
misuse/disuse; the Trust-in-Automation model; XAI.\\
\textbf{The four papers:} Siedl et al.\ (exoskeleton diary study, HLM —
usefulness but not ease of use, attractiveness early, then social feedback
over time); Liu (dependence on generative AI in 3 stages,
usefulness/anthropomorphism $\to$ fear via dependence, literacy buffers); Li
et al.\ (N$=$244, identity + emotion $\to$ empathy $\to$ helping the chatbot);
Li \& Lee (purpose $>$ performance violations; apology \emph{with explanation}
repairs).\\
\textbf{Methods buzzwords:} self-efficacy (belief in your own ability);
psychology's 4 goals (describe/explain/predict/modify); experiment =
manipulation + random assignment + control $\Rightarrow$ causality; jangle
(different names, same thing) vs.\ jingle (same name, different things); the
replicability crisis (p-hacking, publication bias, HARKing, small samples).
\end{notebox}


% ==========================================================
%  CHAPTER 6 — OLD EXAM QUESTIONS
% ==========================================================
\chapter{Old Exam Questions}

A collection of questions from past Robopsychology exams (SS24, the 2023
Moodle exam, and the 2021 online exam), grouped by topic with the
\textbf{correct answer in bold}. \textbf{SS24} maps almost one-to-one onto the
current four-topic structure and is your best predictor. \textbf{2023} adds
more research-methods depth (plus a \emph{creativity} topic that's no longer in
the syllabus). \textbf{2021} comes from an older, general-psychology curriculum
(sleeper effect, eyewitness memory, priming, short-term memory, the mind--body
problem, somatosensation, face processing) — those specific topics are
\textbf{no longer examined}, so only its recurring items (TAM, trust, uncanny
valley, BPN, the randomized experiment) are kept here.

\section{Needs, motivation \& acceptance}

\begin{enumerate}[leftmargin=1.6em, itemsep=4pt]
\item According to Maslow, which is a \textbf{growth} need? — Self-Actualization
/ Love \& Belonging / Safety / Esteem. \hfill\textbf{$\to$ Self-Actualization}
\item Sheldon et al.\ introduced 10 psychological needs. Which is
\textbf{NOT} included? — Pleasure-Stimulation / Physical-Thriving /
Money-Luxury / Identification. \hfill\textbf{$\to$ Money-Luxury}
\item What are the three Basic Psychological Needs (BPN)? — Freedom, Social
Interaction, Achievement / \textbf{Autonomy, Relatedness, Competence} / Survival,
Growth, Self-Actualization / Safety, Love, Esteem.
\item According to SDT, what are two types of motivation? —
Cognitive/affective / Positive/negative / \textbf{Autonomous and controlled} /
Primary/secondary.
\item According to TAM, the two key variables determining adoption are —
\textbf{Perceived Usefulness and Ease of Use} / Cost \& Popularity / Efficiency
\& Safety / Enjoyment \& Autonomy. \emph{(2021 variant: the CORRECT TAM
statement is ``ease of use and perceived usefulness determine acceptance''.)}
\item Which best describes \textbf{self-efficacy}? — \textbf{The belief in one's
ability to perform tasks} / tendency to avoid challenges / job satisfaction /
desire for approval.
\item Self-efficacy \& exoskeletons — which statement is \textbf{incorrect}? —
\textbf{``Exoskeleton use does not impact task-specific self-efficacy at all.''}
(The other three — relief, perceived usefulness, ``not automatic'' — are
correct.)
\item Name of the first robotic exoskeleton (U.S.\ Army + GE, 1960s)? —
\textbf{Hardiman} / Guardian XO / HAL / BLEEX.
\item \emph{(Open, 4 pts — recurs every year)} Name the three BPN and the
theory they come from, then discuss with an AI example how AI can threaten or
help fulfil one of them. \emph{[SDT; competence/autonomy/relatedness; e.g.\
recommender algorithms can threaten autonomy if decisions are opaque, or an
assistive robot can support an elderly person's autonomy.]}
\end{enumerate}

\section{AI literacy}

\begin{enumerate}[leftmargin=1.6em, itemsep=4pt]
\item Research on \textbf{computer anxiety}: people with \emph{less} computer
anxiety are usually — better educated / \textbf{more self-confident} / younger
than 45 / grown up in Asia.
\item The Scandinavian project ``Q'' reduces gender bias in AI assistants by —
more female developers / \textbf{developing a gender-neutral synthetic voice} /
avoiding speech / controlling training data. \emph{(Older curriculum item;
illustrates voice-gender bias.)}
\end{enumerate}

\section{Anthropomorphism \& social AI}

\begin{enumerate}[leftmargin=1.6em, itemsep=4pt]
\item What did participants watch in Heider \& Simmel? —
\textbf{Geometric shapes moving in/around a square} / vehicles on a map /
dancing characters / migrating animals.
\item The two dimensions of mind perception in a machine are \emph{agency} and
\emph{experience}. (T/F) \hfill\textbf{$\to$ True}
\item Which is \textbf{NOT} an assumption of Reeves \& Nass' CASA paradigm? —
people behave toward computers as toward people / users treat machines as
social entities even if they find it inappropriate / social-psychological
phenomena transfer to HMI / \textbf{social perceptions of machines are conscious
and deliberate}. \emph{(CASA reactions are automatic/unconscious.)}
\item \emph{(Open, 4 pts)} Give a short definition of \textbf{anthropomorphism}
and name \textbf{three assumed consequences}, each with a practical example.
\emph{[Attributing human traits to non-humans; consequences: moral care \&
concern; responsibility/blame; social influence.]}
\item \emph{(Open, 4 pts)} Main assumptions of the \textbf{Uncanny Valley}
hypothesis, Mori's practical advice, and what it suggests for AI design.
\emph{[Likability rises with human-likeness up to a point, then drops sharply
into eeriness (movement deepens it); Mori advised \emph{not} to aim for perfect
human likeness; design AI to be clearly distinguishable from humans.]}
\end{enumerate}

\section{Trust in AI}

\begin{enumerate}[leftmargin=1.6em, itemsep=4pt]
\item Per definition, a trust-relevant situation always involves — a shared
goal / at least two persons / \textbf{vulnerability} / anxiety.
\item In human--AI relationships, \textbf{calibrated trust} refers to —
unjustified distrust / marketing-manipulated trust / \textbf{trust proportionate
to the actual trustworthiness of a system} / unwarranted over-trust.
\item Cognition-based trust is characterized by perceiving the trustee as
well-minded, warm-hearted, benevolent, norm-adhering. (T/F) \hfill\textbf{$\to$
False} \emph{(that describes \textbf{affect}-based trust)}.
\item ``Let's talk about trust — which statement is \textbf{FALSE}?'' —
\textbf{``At market launch, with little experience, no individual differences in
trust are expected.''} \emph{(False: propensity to trust drives individual
differences \emph{precisely} when experience is lacking.)} The other options
(vulnerability is needed; affect-based trust can cloud risk assessment;
appearing competent/reliable supports cognitive trust) are true.
\end{enumerate}

\section{Research methods (recurring across all years)}

\begin{enumerate}[leftmargin=1.6em, itemsep=4pt]
\item What is psychology? — \textbf{The study of the human mind and human
behaviour} (not a sub-discipline of medicine/philosophy).
\item The four central goals of scientific psychology are describe, explain,
predict, modify behaviour. (T/F) \hfill\textbf{$\to$ True}
\item Which design investigates cause--effect relationships? — cross-sectional
survey / correlational study / theoretical review / \textbf{randomized
experiment}.
\item A study where half see only a red robot and half only a blue robot is a
\textbf{between-subjects} design. (T/F) \hfill\textbf{$\to$ True}
\item \emph{Representativeness} = how well the \textbf{sample reflects the
population} (single-choice favourite). \emph{Jangle} fallacy = different names,
same construct; \emph{jingle} = same name, different constructs.
\item The \textbf{replicability crisis} and its causes (p-hacking, publication
bias, HARKing, small/WEIRD samples) — recurring open/MC item in 2023.
\item \emph{(Open, 4 pts)} What advantage does the scientific experiment offer
over correlational studies, and how did a named study (e.g.\ the XAI/trust
paper, Leichtmann et al.) meet $\ge$2 characteristics of an experiment?
\emph{[Causal inference; via manipulation of the IV + random assignment +
control/comparison group.]}
\end{enumerate}

\begin{notebox}[title={Highest-yield, confirmed-recurring questions}]
Across SS24 / 2023 / 2021 the same items keep appearing: \textbf{(1)} the
open ``three BPN + SDT + AI example'' question (every year); \textbf{(2)} TAM =
usefulness + ease of use; \textbf{(3)} a trust \textbf{FALSE-statement} /
\textbf{calibrated trust} item; \textbf{(4)} the open anthropomorphism
definition + 3 consequences; \textbf{(5)} Uncanny Valley (Mori);
\textbf{(6)} CASA assumptions / mind perception agency \& experience;
\textbf{(7)} Heider \& Simmel = moving geometric shapes; \textbf{(8)}
between-subjects T/F and randomized experiment = causality. Master these first.
\end{notebox}


\end{document}
